%% AMS-LaTeX Created with the Wolfram Language : www.wolfram.com

\documentclass[12pt]{article}

\usepackage[letterpaper,margin=1in]{geometry}
\usepackage{amsmath, amssymb, graphics, setspace}
\usepackage{graphicx} 
\graphicspath{ {./figs/} }
\usepackage{tcolorbox}
\usepackage{bm}
\usepackage[framed,numbered,autolinebreaks,useliterate]{mcode}
\newcommand{\mathsym}[1]{{}}
\newcommand{\unicode}[1]{{}}

\usepackage{hyperref}
\hypersetup{
    colorlinks=true,
    linkcolor=blue,
    filecolor=magenta,      
    urlcolor=cyan,
    pdftitle={Overleaf Example},
    pdfpagemode=FullScreen,
}

\def\mathbi#1{\textbf{\em #1}}

\begin{document}
\doublespacing	

\title{CE 401/5501 Homework 12\\
\vspace{2in}
\textbf{Name: \rule{2in}{0.15mm}}}
\author{}
\maketitle

\newcommand*{\I}{\imath} % imaginary unit i

\newpage
\doublespacing
\section*{ML for Regression}
Review the \href{https://umsystem.instructure.com/courses/280496/files/34117291?module_item_id=9971137}{Topic 14 class slides} carefully. Answer the following conceptual questions (which are either multiple-choice,  fill-in-the-blank, or short-answer questions).

\subsection*{Q1}
\textbf{(Multiple Choice)} Comparing supervised classification and machine learning for regression, which statements are correct? (Select all that apply)

\begin{itemize}
  \item[(A)] Both map input variables $\bm{x}$ to outputs $y$.
  \item[(B)] Classification outputs are continuous, regression outputs are categorical.
  \item[(C)] Regression outputs are continuous, classification outputs are categorical.
  \item[(D)] Both can use machine learning models like neural networks.
  \item[(E)] Only regression requires minimizing a loss function.
\end{itemize}

\vspace{1em}
\subsection*{Q2}
\textbf{(Short Answer)} In fluid mechanics, the Manning equation is given by:

\[
V = \frac{1}{n} R^{2/3} S^{1/2}
\]
where:
\begin{itemize}
  \item $V$: velocity (m/s)
  \item $R$: hydraulic radius (m)
  \item $S$: slope of the energy grade line (unitless)
  \item $n$: Manning roughness coefficient (unitless)
\end{itemize}

Discuss how this nonlinear equation can be transformed into a linear regression problem using a suitable transformation.

\vspace{1em}
\subsection*{Q3}
\textbf{(Short Answer)} If you have a very complex nonlinear relationship to fit due to the underlying different scenarios (e.g., the bearing capacity of a shallow foundation, for which many scenarios exist, e.g., for a square foundation, a strip foundation, etc.). Under each scenario, there is a design tradition that engineers desire a simple empirical equation. 


In the meantime, only a small dataset is available.

You are choosing between a neural network (MLP) and a support vector machine (SVR), which strategy would you suggest? Briefly explain your reasoning.

\vspace{1em}
\subsection*{Q4}
\textbf{(Fill-in-the-Blank)} A decision tree is a model that splits data based on feature values.

Complete the sentence below:

\begin{quote}
A decision tree is good for problems where interpretability and handling both \underline{\hspace{2cm}} and \underline{\hspace{2cm}} variables are important, and where the underlying relationship between inputs and outputs is \underline{\hspace{2cm}}.
\end{quote}


\newpage
\section*{Deep Neural Networks - Convolutional Neural Networks}


Review class materials on \href{https://umsystem.instructure.com/courses/280496/files/34117289?module_item_id=9971136}{convolutional neural networks (CNNs)}. Answer the following conceptual and calculation questions carefully.


\subsection*{Q1}
\textbf{(Multiple Choice)} Why do Convolutional Neural Networks (CNNs) generally work well for processing images?

\begin{itemize}
  \item[(A)] CNNs use very large fully connected layers.
  \item[(B)] CNNs are inspired by the structure of the visual cortex in the brain, which processes images in a hierarchical manner.
  \item[(C)] CNNs memorize every pixel individually.
  \item[(D)] CNNs are designed to ignore spatial relationships in data.
\end{itemize}

\vspace{1em}
\subsection*{Q2}
\textbf{(Fill-in-the-Blank)}

Besides fully connected layers, CNNs include specialized layers to process images. Fill in the blanks:

\begin{quote}
CNNs commonly use \underline{\hspace{2cm}} layers and \underline{\hspace{2cm}} layers to extract and reduce spatial features.
\end{quote}

\vspace{1em}
\subsection*{Q3}
\textbf{(Calculation)}

\textit{Note: In convolution, the stride refers to how many pixels the filter moves each time. A stride of 1 means the filter `walks' 1 pixel at a time horizontally or vertically.}

Given the following $5 \times 5$ input image matrix:

\[
\begin{bmatrix}
1 & 2 & 3 & 0 & 1 \\
4 & 5 & 6 & 1 & 0 \\
7 & 8 & 9 & 2 & 1 \\
1 & 3 & 5 & 2 & 0 \\
2 & 4 & 6 & 1 & 1 \\
\end{bmatrix}
\]

and the following $3 \times 3$ kernel:

\[
\begin{bmatrix}
1 & 0 & -1 \\
1 & 0 & -1 \\
1 & 0 & -1 \\
\end{bmatrix}
\]

Assume stride = 1, no padding. Calculate the resulting feature map (\textit{hint: it will lead to a $3\times 3$ feature map}).

\vspace{1em}
\subsection*{Q4}
\textbf{(Max Pooling)}

\textit{Note: In max pooling, the stride determines how far the pooling window moves each time.}

Given the following $5 \times 5$ feature map:

\[
\begin{bmatrix}
1 & 3 & 2 & 1 & 5 \\
4 & 6 & 5 & 0 & 2 \\
7 & 2 & 8 & 3 & 1 \\
2 & 9 & 1 & 4 & 7 \\
0 & 6 & 2 & 8 & 3 \\
\end{bmatrix}
\]

Apply $3 \times 3$ max pooling with stride = 1. Write down the resulting pooled feature map.

\end{document}
